\chapter{Parametrii experimentali}\label{SetariExperimentale}
In acest capitol vom descrie in primul rand parametrii tuturor obiectelor din cadrul scenei care pot influenta simularea.
Ulterior vom descrie setul de simulari, scopul fiecareia, si setul de parametri care le definesc.
\section{Main Camera si Walls}
\begin{itemize}
    \item \texttt{Size} : schimba marimea GameView-ului prezentat de camera, si deci schimba si marimea zonei in care se genereaza fragmente de energie si creaturi.
    \item \texttt{Wall transform}: peretii doar trebuie scalati si repozitionati astfel incat sa inchida complet spatiul din interior, dar sa fie putin inafara view-ului camerei principale.
\end{itemize}

\section{FoodManager}
\begin{itemize}
    \item \texttt{Food Prefab}: prefab-ul obiectului ce va fi folosit pe post de mancare
    \item \texttt{Total Food}: numar intreg ce reprezinta cate obiect Food Prefab vor fi generate si pozitionate in scena. Acest numar ramane constant pe parcursul unei simulari.
\end{itemize}

\section{Prey si Predator}
\begin{itemize}
    \item \texttt{Acceleration Factor}: un factor ce scaleaza viteza unei creaturi, care in mod implicit e intre -1 si 1.
    \item \texttt{Turn Factor}: cu cate grade se poate roti o creatura intr-o directie.
    \item \texttt{Smoothing Factor}: factor ce afecteaza fluiditatea miscarilor, atat a rotatiilor cat si a deplasarilor inainte si inapoi.
    \item \texttt{Network}: reteau atribuita creaturii.
    \item \texttt{InputNodes}: numarul de noduri din stratul de intrarea al retelei.
    \item \texttt{OutputNodes}: numarul de noduri din stratul de iesire al retelei.
    \item \texttt{Total Energy}: cantitate de energie atribuita creaturii initial.
    \item \texttt{Energy For Food}: cata energie este obtinuta la coliziunea cu o sursa de mancare.
    \item \texttt{Speed Consumption Multiplier}: un factor ce scaleaza consumul de energie al creaturii in dependenta de viteza cu care aceasta se misca.
    \item \texttt{View Distance}: parametru al senzorului "FoW", stabileste raza din jurul creaturii in care se pot detecta tintele.
    \item \texttt{View Angle}: parametru al "FoW", stabileste unghiul din fata creaturii in care sunt detecatate tintele.
    \item \texttt{Angle Splits}: parametru al "FoW", stabileste in cate zone egale de detectie este impartit \textbf{View Angle}.
    \item \texttt{Fow Targets Mask}: parametru al "FoW", stabileste pe ce Layer va avea loc detectia, adica care sunt tintele senzorului.
    \item \texttt{Wall Ray Distance}: parametru al senzorului de detectie a peretelui, stabileste distanta pe care este proiectat Raycast-ul.
    \item \texttt{Sprite Color}: culoarea sprite-ului atasat creaturii.
\end{itemize}

\section{NeatManager}
\begin{itemize}
    \item \texttt{Creature Prefab}: Prefab-ul din care se vor construi membrii popoulatiei
    \item \texttt{Input Nodes}: numarul de noduri pentru stratul de intrare al retelelor
    \item \texttt{Output Nodes}: numarul de noduri pentru stratul de intrare al retelelor
    \item \texttt{Init Weight Range}: valoare $x$ care defineste intervalul [$-x$, $x$], intre care vor fi initializate ponderile retelelor.
    \item \texttt{Mutation Prob}: probabilitatea ca un genom sa sufere mutatii la aplicare operatorului de mutatie.
    \item \texttt{Weight Perturbance Prob}: probabilitate $p$ a fiecarei gene de conexiune din genom de a suferi o perturbare. Probabilitate ca unei ponderi sa ii fie reatribuite o valoare aleatorie e $1-p$.
    \item \texttt{Perturbance Range}: valoare $x$ care defineste intervalul [$-x$, $x$], intre care se alege o valoare aleatorie pentru perturbarea unei ponderi.
    \item \texttt{Reassignment Range}: valoare $x$ care defineste intervalul [$-x$, $x$], intre care se alege o valoare aleatorie folosita pentru resetarea ponderii.
    \item \texttt{New Node Prob}: probabilitatea ca unui genom sa ii fie aplicat mutatia structurala de adaugare a unui nou nod.
    \item \texttt{New Connection Prob}: probabilitatea ca unui genom sa ii fie aplicat mutatia structurala de adaugare a unei noi conexiuni.
    \item \texttt{Inherid Disabled Gene Prob}: probabilitatea in cadrul crossover, daca una dintre doua gene aliniate "matching" este dezactivat, copilul va mosteni o gena dezactivata.
    \item \texttt{Chance Mutation No Crossover}: probabilitate ca un nou membru din populatie sa nu fie generat prin crossover, ci doar prin operatorul de mutatie aplicat nodului parinte
    \item \texttt{Inter Species Mating Rate}: probabilitate in timpul ca operatii de crossover ca al doilea parinte sa fie ales dintr-o specie aleatorie
    \item \texttt{Species Dropoff Age}: numarul de generatii pe care o specie il poate supravetui daca aceasta nu prezinta imbunatatiri ale fitness-ului celui mai bun membru.
    \item \texttt{Survival Treshold}: procentul din membrii cei mai buni ai unei specii care sunt pastrati pentru crearea genom-urilor populatiei din generatia viitoare.
    \item \texttt{Species Target Count}: Numarul tinta de specii in cadrul populatiei.
    \item \texttt{Compatibility Treshold}: pragul de compatibilitate $\delta_t$. Daca distanta de compatibilitate dintre doua genom-uri este mai mare decat aceasta valoare, atunci genom-urile fac parte din specii diferite.
    \item \texttt{Min Compatibility Treshold}: valoarea minima pentru pragul de compatibilitate
    \item \texttt{Max Compatibility Treshold}: valoarea maxima a pragului de compatibilitate
    \item \texttt{Compatibility Modif}: valoare cu care este modificat pragul de compatibilitate, daca suntem sub numarul tinta de specii, $\delta_t$ va scadea cu atat la fiecare generatie, iar daca suntem peste numarul tinta, $\delta_t$ va creste cu atat.
    \item \texttt{Excess Coef}: valoarea $c1$ din formula \ref{eq:compatibility}
    \item \texttt{Disjoint Coef}: valoarea $c2$ din formula \ref{eq:compatibility}
    \item \texttt{Weight Diff Coef}: valoarea $c3$ din formula \ref{eq:compatibility}
    \item \texttt{Top Members Selection Cutoff}: daca o specie are mai mult de atatia membri, atunci genomul celui mai bun membru al speciei va trece in urmatoarea generatie fara ai fi aplicati operatorii de crossover sau de mutatie.
    \item \texttt{Starting Energy}: cantitatea de energie cu care incepe fiecare membru din populatie.
    \item \texttt{Energy For Food}: cata energie este obtinuta la coliziunea creaturii cu o sursa de mancare.
    \item \texttt{Starting Population}: cate creaturi sunt in populatia initiala a fiecarei generatii.
\end{itemize}